% Generated draft: requires mathematical and source review.
\documentclass[11pt]{article}
\usepackage[margin=25mm]{geometry}
\usepackage{fontspec}
\setmainfont{Latin Modern Roman}
\usepackage{amsmath,amssymb,mathtools,unicode-math}
\setmathfont{Latin Modern Math}
\usepackage{enumitem,needspace,xurl,hyperref}
\setlength{\parindent}{0pt}
\setlength{\parskip}{5pt}
\setlength{\emergencystretch}{3em}
\newcommand{\N}{\mathbb{N}}
\newcommand{\Z}{\mathbb{Z}}
\newcommand{\Q}{\mathbb{Q}}
\newcommand{\R}{\mathbb{R}}
\newcommand{\C}{\mathbb{C}}
\newcommand{\F}{\mathbb{F}}
\newcommand{\A}{\mathbb{A}}
\newcommand{\PP}{\mathbb{P}}
\newcommand{\E}{\mathbb{E}}
\newcommand{\eps}{\varepsilon}
\newcommand{\dd}{\,\mathrm{d}}
\begin{document}
\section*{UnsolvedMath ID: 6700021}
\textbf{Source problem code:} catalogue entry
\subsection{Definitions and mathematical statement}
The supplied problem requests four changes to an unspecified preceding estimate: remove its spin assumption; remove its uniform-positivity assumption; relax completeness; and determine its sharp constant, denoted \(\mathrm{const}_n\), with the indicated dependence on the K-theory class of \(L\).

No inequality is actually supplied: ``[estimate]'' is a placeholder in the statement. The underlying space, the object \(L\), the relevant K-theory group and class, all norms and quantities in the estimate, and the existing hypotheses are absent. Thus even a schematic designation \(C(n,[L])\) cannot be assigned an extremal definition until the inequality and its admissible objects are known. No direction of inequality or normalization is invented.

For context, a spin assumption ordinarily concerns a lift of an oriented tangent bundle to a spin bundle. Uniform positive scalar curvature would mean a bound \(\operatorname{Scal}_g\ge\sigma>0\), stronger than mere pointwise positivity, but the excerpt does not identify which quantity its positivity condition concerns or what replacement is sought. Likewise relaxing completeness requires specified boundary or end assumptions. These contextual meanings do not supply missing source hypotheses.

The summary interprets the background as a Dirac/K-area index estimate and reports special sphere/spin-case progress while saying no unified sharp spin-free result was verified. It gives no exact estimate or original theorem reference. The 2017 open-item label and later open/triage assessment are source information, not an independently verified statement of a complete theorem.
\subsection{Short English statement}
Generalize a preceding index-type estimate by removing spin and uniform positivity, weakening completeness, and finding the sharp constant depending on the relevant K-theory class. The estimate and its objects are missing, so the requested theorem cannot be completed from this record.
\subsection{Sources}
\begin{itemize}
\item Local registry: \url{data/unsolvedmath/problems.json}, permanent UnsolvedMath ID \(6700021\), source code catalogue entry; supplied statement, background, and research summary, updated 2026-08-06.
\item Misha Gromov, \emph{101 Questions, Problems and Conjectures around Scalar Curvature} (2017), printed tag [?21], occurrence 21, PDF page 27: \url{https://www.ihes.fr/~gromov/wp-content/uploads/2018/08/101-problemsOct1-2017.pdf}. Supplied normalized PDF-text extraction accessed 2026-07-29; described as an incomplete source.
\item The summary refers to Dirac/K-area bounds and sphere/spin cases but supplies no exact theorem or separate bibliographic citation.
\end{itemize}
\end{document}
